%%%%%%%%%%%%%%%%%%%%%%%%%%%%%%%%%%%%%%%%%
% Short Sectioned Assignment LaTeX Template Version 1.0 (5/5/12)
% This template has been downloaded from: http://www.LaTeXTemplates.com
% Original author:  Frits Wenneker (http://www.howtotex.com)
% License: CC BY-NC-SA 3.0 (http://creativecommons.org/licenses/by-nc-sa/3.0/)
%%%%%%%%%%%%%%%%%%%%%%%%%%%%%%%%%%%%%%%%%

% \documentclass[paper=a4, fontsize=11pt]{scrartcl} % A4 paper and 11pt font size
\documentclass[11pt, a4paper]{book}
\usepackage[T1]{fontenc} % Use 8-bit encoding that has 256 glyphs
\usepackage[utf8]{inputenc}
\usepackage{fourier} % Use the Adobe Utopia font for the document - comment this line to return to the LaTeX default
\usepackage{listings} % para insertar código con formato similar al editor
\usepackage[spanish, es-tabla]{babel} % Selecciona el español para palabras introducidas automáticamente, p.ej. "septiembre" en la fecha y especifica que se use la palabra Tabla en vez de Cuadro
\usepackage{url} % ,href} %para incluir URLs e hipervínculos dentro del texto (aunque hay que instalar href)
\usepackage{graphics,graphicx, float} %para incluir imágenes y colocarlas
\usepackage[gen]{eurosym} %para incluir el símbolo del euro
\usepackage{cite} %para incluir citas del archivo <nombre>.bib
\usepackage{enumerate}
\usepackage{hyperref}
\usepackage{graphicx}
\usepackage{tabularx}
\usepackage{booktabs}

\usepackage[table,xcdraw]{xcolor}
\hypersetup{
	colorlinks=true,	% false: boxed links; true: colored links
	linkcolor=black,	% color of internal links
	urlcolor=cyan		% color of external links
}
\renewcommand{\familydefault}{\sfdefault}
\usepackage{fancyhdr} % Custom headers and footers
\pagestyle{fancyplain} % Makes all pages in the document conform to the custom headers and footers
\fancyhead[L]{} % Empty left header
\fancyhead[C]{} % Empty center header
\fancyhead[R]{Me Llamo Así} % My name
\fancyfoot[L]{} % Empty left footer
\fancyfoot[C]{} % Empty center footer
\fancyfoot[R]{\thepage} % Page numbering for right footer
%\renewcommand{\headrulewidth}{0pt} % Remove header underlines
\renewcommand{\footrulewidth}{0pt} % Remove footer underlines
\setlength{\headheight}{13.6pt} % Customize the height of the header

\usepackage{titlesec, blindtext, color}
\definecolor{gray75}{gray}{0.75}
\newcommand{\hsp}{\hspace{20pt}}
\titleformat{\chapter}[hang]{\Huge\bfseries}{\thechapter\hsp\textcolor{gray75}{|}\hsp}{0pt}{\Huge\bfseries}
\setcounter{secnumdepth}{4}
\usepackage[Lenny]{fncychap}


\begin{document}

	% Plantilla portada UGR
	\begin{titlepage}
\newlength{\centeroffset}
\setlength{\centeroffset}{-0.5\oddsidemargin}
\addtolength{\centeroffset}{0.5\evensidemargin}
\thispagestyle{empty}

\noindent\hspace*{\centeroffset}\begin{minipage}{\textwidth}

\centering
\includegraphics[width=0.9\textwidth]{logos/logo_ugr.jpg}\\[1.4cm]

\textsc{ \Large TRABAJO FIN DE GRADO\\[0.2cm]}
\textsc{ GRADO EN INGENIERIA INFORMATICA}\\[1cm]

{\Huge\bfseries Título \\}
\noindent\rule[-1ex]{\textwidth}{3pt}\\[3.5ex]
{\large\bfseries Subtítulo }
\end{minipage}

\vspace{2.5cm}
\noindent\hspace*{\centeroffset}
\begin{minipage}{\textwidth}
\centering

\textbf{Autor}\\ {Estudiante}\\[2.5ex]
\textbf{Director}\\ {Tutor(a)(es)}\\[2cm]
\includegraphics[width=0.3\textwidth]{logos/etsiit_logo.png}\\[0.1cm]
\textsc{Escuela Técnica Superior de Ingenierías Informática y de Telecomunicación}\\
\textsc{---}\\
Granada, Junio de 201x
\end{minipage}
\end{titlepage}


	% Plantilla prefacio UGR
	\thispagestyle{empty}

\begin{center}
{\large\bfseries \titulo \\ \subtitulo }\\
\end{center}
\begin{center}
\minombre\\
\end{center}

%\vspace{0.7cm}

\vspace{0.5cm}
\noindent\textbf{Palabras clave}: \textit{galería de arte, aplicación web, React, NodeJS, MySQL, Docker}
\vspace{0.7cm}

\noindent\textbf{Resumen}

\bigskip
El objetivo de este trabajo es el desarrollo de una aplicación web
que sirva como una galería de arte online en la que los artistas podrán publicar
sus obras para darlas a conocer. Los usuarios podrán ver y valorar las obras para poder
dar feedback a los artistas. El propósito de esta aplicación es que los artistas puedan
dar a conocer sus obras y que galerías de arte convencionales puedan descubrir nuevos
artistas para exponer sus obras.

\cleardoublepage

\begin{center}
	{\large\bfseries \titulo \\ \subtituloingles }\\
\end{center}
\begin{center}
	\minombre\\
\end{center}
\vspace{0.5cm}
\noindent\textbf{Keywords}: \textit{art gallery, web app, React, NodeJS, MySQL, Docker}
\vspace{0.7cm}

\noindent\textbf{Abstract}

\bigskip
The objective of this work is the development of a web application
that will serve as an online art gallery that serves as an online art
gallery where artists can publish their works to make them known.
Users will be able to view and rate the artworks in order to give
feedback to the artists. The purpose of this application is for artists
to be able to to make their works known and for conventional art
galleries to discover new artists to exhibit their works.


\cleardoublepage

\thispagestyle{empty}

\noindent\rule[-1ex]{\textwidth}{2pt}\\[4.5ex]

D. \textbf{\tutor}, catedrático del departamento de Ciencias de la
Computación e Inteligencia Artificial de la Universidad de Granada.

\vspace{0.5cm}

\textbf{Informo:}

\vspace{0.5cm}

Que el presente trabajo, titulado \textit{\textbf{\titulo}},
ha sido realizado bajo mi supervisión por \textbf{\minombre}, y autorizo la defensa de dicho trabajo ante el tribunal
que corresponda.

\vspace{0.5cm}

Y para que conste, expiden y firman el presente informe en Granada a Junio de 2023.

\vspace{1cm}

\textbf{El director: }

\vspace{5cm}

\noindent \textbf{\tutor}

\chapter*{Agradecimientos}






	% Índice de contenidos
	\newpage
	\tableofcontents

	% Índice de imágenes y tablas
	\newpage
	\listoffigures

	% Si hay suficientes se incluirá dicho índice
	\listoftables 
	\newpage

	% Introducción 
	\chapter{Introducción}

La tecnología está presente en casi todos los aspectos de nuestra vida a día de hoy. De hecho,
es un recurso que permite conectarnos con realidades que, de otra manera, quedarían
invisibilizadas. Sin embargo, en el mundo del arte nos encontramos con una circunstancia algo
diferente, ya que no se suele promocionar en los medios de comunicación a los que accede la
gran mayoría de la población. En este sentido, aunque entre sus amantes sí que existe un
mayor interés por seguir conociendo nuevos artistas, quizá para el resto de la población
general queda al margen. De esta manera, si una persona no tiene la iniciativa suficiente
como para conectarse día a día a las novedades, puede resultarle más complejo mantenerse
actualizado. 
 
\section{Motivación}  
 
En este contexto, surge la idea de desarrollar una aplicación web que permita a los artistas
publicar sus obras y promocionar su trabajo. Así, cualquier persona puede tener un fácil
acceso a las mismas, aumentando así la probabilidad de que más población se acerque a este
mundo. Además, para todos aquellos amantes del arte también es una forma de, en definitiva,
mejorar y reforzar la cercanía con los autores. Por ende, la principal motivación que inicia
este proyecto es crear una plataforma que permita promocionar las obras de pequeños artistas,
que sea intuitiva y fácil de usar para cualquier persona (ya sea quien enseña su trabajo o
quien disfruta conociéndolo). En definitiva, lo que se desea es acercar la cultura a la
población y viceversa, a través del uso de la tecnología.  
 
\section{Descripción del problema} 
 
Los artistas independientes suelen tener dificultades para darse a conocer y las herramientas
que tienen a su disposición para ello no son específicas para este fin. La mayoría de
artistas utilizan redes sociales como \textit{Instagram} \cite{instagram}, \textit{Facebook}
\cite{facebook} o \textit{Twitter} \cite{twitter} para publicar sus obras. Sin embargo, estas
aplicaciones no están diseñadas con este propósito; la mayoría de usuarios no las utilizan
para descubrir nuevos artistas sino para otros fines. Así, es difícil encontrar obras entre
la gran cantidad de contenido que se publica en estas redes.  
 
Por otro lado, las galerías de arte convencionales o estudios fotográficos no cuentan con
herramientas para poder descubrir nuevos artistas en sus exposiciones. En la mayoría de casos,
los artistas que muestran su trabajo en estos lugares ya tienen cierto reconocimiento y se
han dado a conocer anteriormente en otras galerías. Sin embargo, aquellos que se están
iniciando en este mundo y que quizá no han tenido las mismas oportunidades y/o facilidades
que sus colegas para acceder a estos recursos, a veces simplemente por mera cuestión de
suerte, tienen más dificultades para empezar a trabajar de lo que tanto les gusta.  

Además, los grandes centros urbanos suelen ser los lugares que alojan eventos culturales para
promocionar el arte independiente. Sin embargo, a día de hoy es más difícil encontrar esta
misma situación en ciudades y pueblos más pequeños. En ellos, el único medio que tienen las
personas para conocer nuevos artistas es a través de internet, pero, como se ha descrito
anteriormente, no existe ninguna plataforma específica para este fin.

\section{Objetivos}

El propósito de este proyecto es el desarrollo de una aplicación web realizando todo
el proceso de ingeniería del software. Se realizará la fase de planificación, análisis,
diseño e implementación de la aplicación.

También se prentende aprender a utilizar tecnologías para el desarrollo de aplicaciones
web, enfocando el proyecto a un desarrollo full-stack. Se utilizarán tecnologías de las
que no se tenga mucho conocimiento y se quiera aprender como \textit{React} \cite{react}
o \textit{NodeJS} \cite{nodejs}.

Además de aprender a utilizar estas tecnologías, se pretende aprender a utilizar
de una forma más avanzada \textit{Docker} \cite{docker} y \textit{docker-compose}
\cite{docker-compose} para conectar las distintas partes de la aplicación.

Por último, se pretende desarrollar la aplicación de manera que sea intuitiva y fácil
de usar por personas que no posean grandes conocimientos informáticos. La mayoría de
las aplicaciones web que existen están enfocadas a un público con conocimientos o son
poco intuitivas para el usuario final. Para ello será importante realizar un buen
diseño de la interfaz de usuario que sea atractivo y fácil de usar.

\section{Revisión del estado del Arte}

{ \setlength{\parskip}{7mm} % Customize the space between paragraphs
Existen varias aplicaciones web que sirven como galerías de arte online, pero la mayoría
de ellas están enfocadas a la venta de obras de arte y no tanto a la exposición de obras
y promoción de las mismas.

Un ejemplo de esto es la aplicación web \textit{Artfinder} \cite{artfinder}, que permite
a los artistas vender sus obras a través de la plataforma. Sin embargo, no permite a los
usuarios valorar las obras ni dar feedback a los artistas, tampoco poder encontrar las obras
con mayor puntuación.

Otra aplicación es \textit{Singulart} \cite{singulart}, que también permite a los artistas
subir sus obras, aunque para ello, primero deben postular y luego ser aceptados por el equipo
de la plataforma. Además, la página advierte a los artistas que el proceso de aceptación
es muy largo ya que reciben demasiadas solicitudes. Por otro lado, los usuarios tampoco
pueden valorar las obras ni dejar sus comentarios. La página actúa como un \textit{e-commerce}
o tienda online, en la que los artistas ponen a la venta sus obras y los usuarios pueden
buscar obras y comprarlas.

En cuanto a las aplicaciones web que permiten a los usuarios valorar las obras, se encuentra
\textit{ArtStation} \cite{artstation}, que es una plataforma para que los artistas suban
sus obras y los usuarios puedan valorarlas. Sin embargo, esta plataforma está enfocada
a artistas digitales solamente. Además, la interfaz está muy sobrecargada y no es intuitiva
para los usuarios.

En definitiva, no hay una gran oferta de aplicaciones web que permitan a los artistas
independientes publicar sus obras para darse a conocer y a los usuarios valorarlas. La
mayoría de aplicaciones web que existen están enfocadas a la venta de obras de arte y no
tanto a la exposición de obras y promoción de los artistas.
}

	% Descripción del problema y hasta donde se llega
	\chapter{Descripción del problema}



	% Estado del arte
	% 	1. Crítica al estado del arte
	% 	2. Propuesta
	\chapter{Estado del arte}

El software libre y sus licencias \cite{gplv3} ha permitido llevar a cabo una expansión del
aprendizaje de la informática sin precedentes.

	
	\chapter{Planificación}

\section{Metodología utilizada}


\section{Temporización}

\section{Seguimiento del desarrollo}


	% Análisis del problema
	% 1. Análisis de requisitos
	% 2. Análisis de las soluciones
	% 3. Solucion propuesta
	% 4. Análisis de seguridad
	\chapter{Análisis del problema}
 


	% Desarrollo bajo sprints: 
	% 	1. Permitir registros y login de usuarios
	% 	2. Desarrollo del sistema de incidencias
	% 	3. Desarrollo del sistema de denuncias administrativas y accidentes
	% 	4. Desarrollo del sistema de croquis
	%   5. Instalación de la aplicación de manera automática
	\chapter{Implementación}

La implementación del software se ha dividido en hitos. Estos, han sido definidos en Github
y cada uno de ellos contiene un grupo de \textit{issues} que se corresponden con las distintas
mejoras que se han ido incorporando al software a lo largo de su desarrollo.\\



	% Presupuesto

	% Conclusiones
	\chapter{Conclusiones y trabajos futuros}



	% Trabajos futuros


	
	\newpage
	\bibliography{bibliografia}
	\bibliographystyle{plain}
	
\end{document}

